\begin{scanpage}{0237}{219}
\begin{center}REFERENCES\end{center}
{\cyrillicfont\small
\noindent 1. Фок В. А. Обобщение отражательных формул на случай отражения произвольной волны от поверхности произвольной формы. ЖЭТФ, 20, 961, 1950.\par
2. Антенны сантиметровых волн, ч. I. Пер. с англ. под ред. Я. Н. Фельда. Изд-во «Советское радио», 1950.\par
3. Ментцер Дж. Р. Дифракция и рассеяние радиоволн. Пер. с англ. под ред. Л. А. Вайнштейна. Изд-во «Советское радио», 1958.\par
4. Вайнштейн Л. А. Электромагнитные волны. Изд-во «Советское радио», 1957.\par
5. Потехин А. И. Некоторые задачи дифракции электромагнитных волн. Изд-во «Советское радио», 1948.\par
6. Baker B. B. and Copson E. T. The mathematical theory of Huygens' principle. Oxford, 1939.\par
7. Уфимцев П. Я. Приближенный расчет дифракции плоских электромагнитных волн на некоторых металлических телах, ч. I. Дифракция на клине и ленте. ЖТФ, 27, № 8, 1840--1849, 1957.\par
8. Уфимцев П. Я. Приближенный расчет дифракции плоских электромагнитных волн на некоторых металлических телах. ч. II. Дифракция на диске и конечном цилиндре. ЖТФ, 28, № 11, 2601--2616, 1958.\par
9. Уфимцев П. Я. Вторичная дифракция на ленте. ЖТФ, 28, № 3, 569--582, 1958.\par
10. Уфимцев П. Я. Вторичная дифракция на диске. ЖТФ, 28, № 3, 583--591, 1958.\par
11. Уфимцев П. Я. Симметричное облучение конечных тел вращения. «Радиотехника и электроника», 6, № 4, 559--567, 1961.\par
12. Майзельс Е. Н., Уфимцев П. Я. Отражение электромагнитных волн круговой поляризации от металлических тел. «Радиотехника и электроника», 5, № 12, 1925--1928, 1960.\par
13. Уфимцев П. Я. Об отражении металлическими телами радиоволн круговой поляризации. «Радиотехника и электроника», № 12, 2094, 1961.\par
14. Уфимцев П. Я. Дифракция плоских электромагнитных волн на тонком цилиндрическом проводнике. «Радиотехника и электроника», 7, № 2, 260--269, 1962.\par
15. Schwarzschild K. Beugung und Polarisation des Lichts. Mathematische Annalen, 55, 177, 1902.\par
16. Франк Ф. и Мизес Р. Дифференциальные и интегральные уравнения математической физики. ОНТИ, 1937.\par
17. Фок В. А. Распределение токов, возбуждаемых плоской волной на поверхности проводника. ЖЭТФ, 15, № 12, 693--702, 1945.\par
18. Фок В. А. Дифракция радиоволн вокруг земной поверхности. М.-Л. Изд-во АН СССР, 1946.\par
19. Вайнштейн Л. А. Дифракция электромагнитных и звуковых волн на открытом конце волновода. Изд-во «Советское радио», 1953.\par
20. Oberhettinger F. On the diffraction and reflection on waves and pulses by wedges and corners. Journ. Research NBS, 61, № 5, 343--365, 1958.\par
21. Ватсон Г. Н. Теория бесселевых функций. Часть первая. И. Л., 1949.\par
22. Pauli W. On asymptotic series for functions in the theory of diffraction of light. Physical Review, 54, № 11, 924--931, 1938.\par
23. Мак-Лахлан Н. В. Теория и приложение функций Матье. И. Л., 1953.\par
24. Meixner J. und Andrejewski W. Strenge Theorie der Beugung ebener elektromagnetischer Wellen an der vollkommen leitenden Kreisscheibe und an der kreisförmigen Öffnung im vollkommen leitenden ebenen Schirm. Annalen der Physik, 7, № 3--4, 157, 1950.\par
25. Andrejewski W. Die Beugung elektromagnetischer Wellen an der leitenden Kreisscheibe und an der kreisförmigen Öffnung im leitenden ebenen Schirm. Zeits. für angewandte Physik, 5, № 5, 178, 1953.\par}
\end{scanpage}
